\documentclass[a4paper,10pt]{article}
\usepackage[utf8]{inputenc}
\usepackage[T1]{fontenc}
\usepackage[spanish]{babel}
\usepackage{geometry}
\usepackage{array}
\usepackage{colortbl}
\usepackage{xcolor}
\usepackage{booktabs}
\usepackage{longtable}
\usepackage{fancyhdr}
\usepackage{tcolorbox}
\usepackage{enumitem}
\usepackage{multicol}
\usepackage{tabularx}

\geometry{left=1.5cm, right=1.5cm, top=2cm, bottom=1.8cm}

% Colores principales por sección
\definecolor{azulverbo}{HTML}{2980B9}
\definecolor{tealconj}{HTML}{16A085}
\definecolor{moradopos}{HTML}{8E44AD}
\definecolor{naranjainterr}{HTML}{E67E22}

% Colores auxiliares
\definecolor{gris}{HTML}{F2F2F2}
\definecolor{grisclaro}{HTML}{F8F8F8}
\definecolor{rojoirr}{HTML}{C0392B}
\definecolor{verdeosc}{HTML}{27AE60}
\definecolor{azulg}{HTML}{3498DB}

\tcbuselibrary{skins, breakable}

% Estilo de caja para trucos/estrategias
\newtcolorbox{trucoverbo}{
  colback=azulverbo!8,
  colframe=azulverbo!70,
  coltitle=white,
  fonttitle=\bfseries,
  title={Truco / Estrategia},
  rounded corners,
  boxrule=0.8pt,
  left=4pt, right=4pt, top=4pt, bottom=4pt,
  before skip=6pt, after skip=6pt
}

\newtcolorbox{trucoconjugacion}{
  colback=tealconj!8,
  colframe=tealconj!70,
  coltitle=white,
  fonttitle=\bfseries,
  title={Truco / Estrategia},
  rounded corners,
  boxrule=0.8pt,
  left=4pt, right=4pt, top=4pt, bottom=4pt,
  before skip=6pt, after skip=6pt
}

\newtcolorbox{trucoposesivo}{
  colback=moradopos!8,
  colframe=moradopos!70,
  coltitle=white,
  fonttitle=\bfseries,
  title={Truco / Estrategia},
  rounded corners,
  boxrule=0.8pt,
  left=4pt, right=4pt, top=4pt, bottom=4pt,
  before skip=6pt, after skip=6pt
}

\newtcolorbox{trucointerr}{
  colback=naranjainterr!8,
  colframe=naranjainterr!70,
  coltitle=white,
  fonttitle=\bfseries,
  title={Truco / Estrategia},
  rounded corners,
  boxrule=0.8pt,
  left=4pt, right=4pt, top=4pt, bottom=4pt,
  before skip=6pt, after skip=6pt
}

% Encabezado de tarjeta de verbo (media página aprox.)
\newcommand{\verboheader}[3]{%
  % #1 = número, #2 = verbo en mayúsculas, #3 = subtítulo (conjugación, tipo)
  \noindent
  \colorbox{azulverbo!90}{\parbox{\dimexpr\linewidth-2\fboxsep}{%
    \centering\color{white}\bfseries\large Tarjeta #1 --- #2\\[2pt]
    \normalsize #3%
  }}
  \vspace{4pt}
}

\newcommand{\conjheader}[2]{%
  \noindent
  \colorbox{tealconj!90}{\parbox{\dimexpr\linewidth-2\fboxsep}{%
    \centering\color{white}\bfseries\large Tarjeta conjugación #1\\[2pt]
    \normalsize #2%
  }}
  \vspace{4pt}
}

\newcommand{\posheader}[2]{%
  \noindent
  \colorbox{moradopos!90}{\parbox{\dimexpr\linewidth-2\fboxsep}{%
    \centering\color{white}\bfseries\large Tarjeta posesivos #1\\[2pt]
    \normalsize #2%
  }}
  \vspace{4pt}
}

\newcommand{\interrheader}[2]{%
  \noindent
  \colorbox{naranjainterr!90}{\parbox{\dimexpr\linewidth-2\fboxsep}{%
    \centering\color{white}\bfseries\large Tarjeta interrogativos #1\\[2pt]
    \normalsize #2%
  }}
  \vspace{4pt}
}

% Línea de separación entre 2 tarjetas en la misma página
\newcommand{\tarjetasep}{%
  \vspace{6pt}
  \noindent\rule{\linewidth}{0.4pt}
  \vspace{6pt}
}

\pagestyle{fancy}
\fancyhf{}
\fancyhead[L]{\small\textbf{Nuevo Compañeros 1 --- U03: La Familia}}
\fancyhead[R]{\small\textbf{Tarjetas de gramática --- Caja 3}}
\fancyfoot[C]{\thepage}

\setlength{\parindent}{0pt}
\renewcommand{\arraystretch}{1.2}

\begin{document}

%% ============================================================
%% PORTADA
%% ============================================================

\thispagestyle{empty}
\vspace*{3cm}

\begin{center}
{\Huge\bfseries\textcolor{azulverbo}{GRAMATIPS}}\\[0.6cm]
{\LARGE\bfseries UNIDAD 3: LA FAMILIA}\\[0.4cm]
{\large Caja 3 --- Estación de servicio}\\[0.3cm]
{\normalsize Gramática (p.\,36--37) --- Nuevo Compañeros 1}
\end{center}

\vspace{1.5cm}

\begin{center}
\begin{tabular}{|>{\bfseries\centering\arraybackslash}p{4cm}|c|>{\centering\arraybackslash}p{3cm}|>{\centering\arraybackslash}p{3cm}|}
\hline
\rowcolor{gris}
\textbf{Sección} & \textbf{Tarjetas} & \textbf{Color} & \textbf{Páginas} \\
\hline
\textcolor{azulverbo}{A. Verbos} & 10 & \textcolor{azulverbo}{azul} & 2--6 \\
\hline
\textcolor{tealconj}{A-bis. Conjugación} & 3 & \textcolor{tealconj}{verde azulado} & 7--9 \\
\hline
\textcolor{moradopos}{B. Posesivos} & 3 & \textcolor{moradopos}{morado} & 10--12 \\
\hline
\textcolor{naranjainterr}{C. Interrogativos} & 3 & \textcolor{naranjainterr}{naranja} & 13--15 \\
\hline
\rowcolor{gris}
\textbf{TOTAL} & \textbf{19} & & \\
\hline
\end{tabular}
\end{center}

\vspace{1.5cm}

\begin{center}
\begin{tcolorbox}[colback=gris, colframe=black!40, width=12cm, boxrule=0.5pt]
\centering
\textbf{Conjugación de ESTAR} (irregular --- incluida en cada tarjeta de verbo):\\[4pt]
\begin{tabular}{ll}
yo & estoy \\
tú & estás \\
él / ella / Ud. & está \\
nosotros/-as & estamos \\
vosotros/-as & estáis \\
ellos / ellas / Uds. & están \\
\end{tabular}
\end{tcolorbox}
\end{center}

\newpage

%% ============================================================
%% A. TARJETAS DE VERBOS (azul) — 10 tarjetas, 2 por página
%% ============================================================

%%---------- TARJETA 1: ESTUDIAR ----------
\verboheader{1}{ESTUDIAR}{-ar, 1.\textsuperscript{a} conjugación, regular}

\begin{minipage}[t]{0.48\linewidth}
\small
\begin{tabular}{|>{\bfseries}l|l|}
\hline
\rowcolor{azulverbo!15} Forma & \\
\hline
Infinitivo & estudiar \\
\hline
\rowcolor{grisclaro}
Participio & estudiado \\
\hline
Gerundio & estudiando \\
\hline
\rowcolor{grisclaro}
Estar + ger. & estoy\,/\,estás\,/\,está\,/\ldots\;+ estudiando \\
\hline
\end{tabular}
\end{minipage}
\hfill
\begin{minipage}[t]{0.48\linewidth}
\small
\begin{tabular}{|l|l|}
\hline
\rowcolor{azulverbo!15} \textbf{Presente} & \\
\hline
yo & estudio \\
\hline
\rowcolor{grisclaro}
tú & estudias \\
\hline
él / ella / Ud. & estudia \\
\hline
\rowcolor{grisclaro}
nosotros/-as & estudiamos \\
\hline
vosotros/-as & estudiáis \\
\hline
\rowcolor{grisclaro}
ellos / ellas / Uds. & estudian \\
\hline
\end{tabular}
\end{minipage}

\vspace{4pt}
\small
\textbf{Frases:}
\begin{itemize}[nosep, leftmargin=1em]
\item Presente: \textit{Jorge \textbf{estudia} música en el instituto.}
\item Infinitivo: \textit{A Jorge le gusta \textbf{estudiar} música.}
\item Participio: \textit{Jorge ha \textbf{estudiado} mucho este curso.}
\item Gerundio: \textit{Jorge está \textbf{estudiando} para el examen.}
\end{itemize}

\begin{trucoverbo}
\textbf{Regularidad:} Verbo regular de la 1.\textsuperscript{a} conjugación (-ar). Se forma quitando -ar y añadiendo las terminaciones: \textbf{-o, -as, -a, -amos, -áis, -an}. Estas terminaciones son iguales para TODOS los verbos en -ar. Participio regular (-ado). Gerundio regular (-ando).
\end{trucoverbo}

\tarjetasep

%%---------- TARJETA 2: TRABAJAR ----------
\verboheader{2}{TRABAJAR}{-ar, 1.\textsuperscript{a} conjugación, regular}

\begin{minipage}[t]{0.48\linewidth}
\small
\begin{tabular}{|>{\bfseries}l|l|}
\hline
\rowcolor{azulverbo!15} Forma & \\
\hline
Infinitivo & trabajar \\
\hline
\rowcolor{grisclaro}
Participio & trabajado \\
\hline
Gerundio & trabajando \\
\hline
\rowcolor{grisclaro}
Estar + ger. & estoy\,/\,estás\,/\,está\,/\ldots\;+ trabajando \\
\hline
\end{tabular}
\end{minipage}
\hfill
\begin{minipage}[t]{0.48\linewidth}
\small
\begin{tabular}{|l|l|}
\hline
\rowcolor{azulverbo!15} \textbf{Presente} & \\
\hline
yo & trabajo \\
\hline
\rowcolor{grisclaro}
tú & trabajas \\
\hline
él / ella / Ud. & trabaja \\
\hline
\rowcolor{grisclaro}
nosotros/-as & trabajamos \\
\hline
vosotros/-as & trabajáis \\
\hline
\rowcolor{grisclaro}
ellos / ellas / Uds. & trabajan \\
\hline
\end{tabular}
\end{minipage}

\vspace{4pt}
\small
\textbf{Frases:}
\begin{itemize}[nosep, leftmargin=1em]
\item Presente: \textit{Mi madre \textbf{trabaja} en un hospital.}
\item Infinitivo: \textit{Mi madre tiene que \textbf{trabajar} los sábados.}
\item Participio: \textit{Mi madre ha \textbf{trabajado} muchos años en el hospital.}
\item Gerundio: \textit{Mi padre está \textbf{trabajando} ahora.}
\end{itemize}

\begin{trucoverbo}
\textbf{Regularidad:} Verbo regular de la 1.\textsuperscript{a} conjugación (-ar). Mismas terminaciones que \textit{estudiar}: \textbf{-o, -as, -a, -amos, -áis, -an}. Participio regular (-ado). Gerundio regular (-ando). Los verbos en -ar son los más numerosos del español.
\end{trucoverbo}

\newpage

%%---------- TARJETA 3: HABLAR ----------
\verboheader{3}{HABLAR}{-ar, 1.\textsuperscript{a} conjugación, regular}

\begin{minipage}[t]{0.48\linewidth}
\small
\begin{tabular}{|>{\bfseries}l|l|}
\hline
\rowcolor{azulverbo!15} Forma & \\
\hline
Infinitivo & hablar \\
\hline
\rowcolor{grisclaro}
Participio & hablado \\
\hline
Gerundio & hablando \\
\hline
\rowcolor{grisclaro}
Estar + ger. & estoy\,/\,estás\,/\,está\,/\ldots\;+ hablando \\
\hline
\end{tabular}
\end{minipage}
\hfill
\begin{minipage}[t]{0.48\linewidth}
\small
\begin{tabular}{|l|l|}
\hline
\rowcolor{azulverbo!15} \textbf{Presente} & \\
\hline
yo & hablo \\
\hline
\rowcolor{grisclaro}
tú & hablas \\
\hline
él / ella / Ud. & habla \\
\hline
\rowcolor{grisclaro}
nosotros/-as & hablamos \\
\hline
vosotros/-as & habláis \\
\hline
\rowcolor{grisclaro}
ellos / ellas / Uds. & hablan \\
\hline
\end{tabular}
\end{minipage}

\vspace{4pt}
\small
\textbf{Frases:}
\begin{itemize}[nosep, leftmargin=1em]
\item Presente: \textit{María y yo \textbf{hablamos} mucho por teléfono.}
\item Infinitivo: \textit{Me gusta \textbf{hablar} español con mis amigos.}
\item Participio: \textit{Hoy he \textbf{hablado} con mi abuela.}
\item Gerundio: \textit{María está \textbf{hablando} por teléfono.}
\end{itemize}

\begin{trucoverbo}
\textbf{Regularidad:} Verbo regular de la 1.\textsuperscript{a} conjugación (-ar). Mismas terminaciones que \textit{estudiar} y \textit{trabajar}. Participio regular (-ado). Gerundio regular (-ando). Tres verbos -ar con el mismo patrón confirman la regla: si conoces las terminaciones de uno, conoces las de todos.
\end{trucoverbo}

\tarjetasep

%%---------- TARJETA 4: COMER ----------
\verboheader{4}{COMER}{-er, 2.\textsuperscript{a} conjugación, regular}

\begin{minipage}[t]{0.48\linewidth}
\small
\begin{tabular}{|>{\bfseries}l|l|}
\hline
\rowcolor{azulverbo!15} Forma & \\
\hline
Infinitivo & comer \\
\hline
\rowcolor{grisclaro}
Participio & comido \\
\hline
Gerundio & comiendo \\
\hline
\rowcolor{grisclaro}
Estar + ger. & estoy\,/\,estás\,/\,está\,/\ldots\;+ comiendo \\
\hline
\end{tabular}
\end{minipage}
\hfill
\begin{minipage}[t]{0.48\linewidth}
\small
\begin{tabular}{|l|l|}
\hline
\rowcolor{azulverbo!15} \textbf{Presente} & \\
\hline
yo & como \\
\hline
\rowcolor{grisclaro}
tú & comes \\
\hline
él / ella / Ud. & come \\
\hline
\rowcolor{grisclaro}
nosotros/-as & comemos \\
\hline
vosotros/-as & coméis \\
\hline
\rowcolor{grisclaro}
ellos / ellas / Uds. & comen \\
\hline
\end{tabular}
\end{minipage}

\vspace{4pt}
\small
\textbf{Frases:}
\begin{itemize}[nosep, leftmargin=1em]
\item Presente: \textit{Yo \textbf{como} a las dos y veinte.}
\item Infinitivo: \textit{Vamos a \textbf{comer} en el instituto.}
\item Participio: \textit{Los niños ya han \textbf{comido}.}
\item Gerundio: \textit{Los niños están \textbf{comiendo} en el instituto.}
\end{itemize}

\begin{trucoverbo}
\textbf{Regularidad:} Verbo regular de la 2.\textsuperscript{a} conjugación (-er). Se forma quitando -er y añadiendo: \textbf{-o, -es, -e, -emos, -éis, -en}. Participio regular (-ido). Gerundio regular (-iendo). Comparar con -ar: las terminaciones son parecidas pero con vocal \textit{e} en vez de \textit{a} (excepto \textit{yo}, que siempre es -o).
\end{trucoverbo}

\newpage

%%---------- TARJETA 5: BEBER ----------
\verboheader{5}{BEBER}{-er, 2.\textsuperscript{a} conjugación, regular}

\begin{minipage}[t]{0.48\linewidth}
\small
\begin{tabular}{|>{\bfseries}l|l|}
\hline
\rowcolor{azulverbo!15} Forma & \\
\hline
Infinitivo & beber \\
\hline
\rowcolor{grisclaro}
Participio & bebido \\
\hline
Gerundio & bebiendo \\
\hline
\rowcolor{grisclaro}
Estar + ger. & estoy\,/\,estás\,/\,está\,/\ldots\;+ bebiendo \\
\hline
\end{tabular}
\end{minipage}
\hfill
\begin{minipage}[t]{0.48\linewidth}
\small
\begin{tabular}{|l|l|}
\hline
\rowcolor{azulverbo!15} \textbf{Presente} & \\
\hline
yo & bebo \\
\hline
\rowcolor{grisclaro}
tú & bebes \\
\hline
él / ella / Ud. & bebe \\
\hline
\rowcolor{grisclaro}
nosotros/-as & bebemos \\
\hline
vosotros/-as & bebéis \\
\hline
\rowcolor{grisclaro}
ellos / ellas / Uds. & beben \\
\hline
\end{tabular}
\end{minipage}

\vspace{4pt}
\small
\textbf{Frases:}
\begin{itemize}[nosep, leftmargin=1em]
\item Presente: \textit{¿Cuánta agua \textbf{bebes} en un día?}
\item Infinitivo: \textit{Tienes que \textbf{beber} más agua.}
\item Participio: \textit{Mi abuelo ha \textbf{bebido} mucho café hoy.}
\item Gerundio: \textit{Mi abuelo está \textbf{bebiendo} café.}
\end{itemize}

\begin{trucoverbo}
\textbf{Regularidad:} Verbo regular de la 2.\textsuperscript{a} conjugación (-er). Mismas terminaciones que \textit{comer}: \textbf{-o, -es, -e, -emos, -éis, -en}. Participio regular (-ido). Gerundio regular (-iendo).
\end{trucoverbo}

\tarjetasep

%%---------- TARJETA 6: VIVIR ----------
\verboheader{6}{VIVIR}{-ir, 3.\textsuperscript{a} conjugación, regular}

\begin{minipage}[t]{0.48\linewidth}
\small
\begin{tabular}{|>{\bfseries}l|l|}
\hline
\rowcolor{azulverbo!15} Forma & \\
\hline
Infinitivo & vivir \\
\hline
\rowcolor{grisclaro}
Participio & vivido \\
\hline
Gerundio & viviendo \\
\hline
\rowcolor{grisclaro}
Estar + ger. & estoy\,/\,estás\,/\,está\,/\ldots\;+ viviendo \\
\hline
\end{tabular}
\end{minipage}
\hfill
\begin{minipage}[t]{0.48\linewidth}
\small
\begin{tabular}{|l|l|}
\hline
\rowcolor{azulverbo!15} \textbf{Presente} & \\
\hline
yo & vivo \\
\hline
\rowcolor{grisclaro}
tú & vives \\
\hline
él / ella / Ud. & vive \\
\hline
\rowcolor{grisclaro}
nosotros/-as & vivimos \\
\hline
vosotros/-as & vivís \\
\hline
\rowcolor{grisclaro}
ellos / ellas / Uds. & viven \\
\hline
\end{tabular}
\end{minipage}

\vspace{4pt}
\small
\textbf{Frases:}
\begin{itemize}[nosep, leftmargin=1em]
\item Presente: \textit{Nosotros \textbf{vivimos} en Lima.}
\item Infinitivo: \textit{Mis primos quieren \textbf{vivir} en Madrid.}
\item Participio: \textit{Mis primos han \textbf{vivido} en Arévalo muchos años.}
\item Gerundio: \textit{Mis primos están \textbf{viviendo} en Arévalo.}
\end{itemize}

\begin{trucoverbo}
\textbf{Regularidad:} Verbo regular de la 3.\textsuperscript{a} conjugación (-ir). Terminaciones: \textbf{-o, -es, -e, -imos, -ís, -en}. Participio regular (-ido). Gerundio regular (-iendo). Comparar con -er: son CASI idénticas. La única diferencia está en \textit{nosotros} (-imos, no -emos) y \textit{vosotros} (-ís, no -éis). En las demás personas: iguales.
\end{trucoverbo}

\newpage

%%---------- TARJETA 7: ESCRIBIR ----------
\verboheader{7}{ESCRIBIR}{-ir, 3.\textsuperscript{a} conjugación, regular (participio irregular)}

\begin{minipage}[t]{0.48\linewidth}
\small
\begin{tabular}{|>{\bfseries}l|l|}
\hline
\rowcolor{azulverbo!15} Forma & \\
\hline
Infinitivo & escribir \\
\hline
\rowcolor{grisclaro}
Participio & \textcolor{rojoirr}{\textbf{escrito}} {\footnotesize (IRREGULAR)} \\
\hline
Gerundio & escribiendo \\
\hline
\rowcolor{grisclaro}
Estar + ger. & estoy\,/\,estás\,/\,está\,/\ldots\;+ escribiendo \\
\hline
\end{tabular}
\end{minipage}
\hfill
\begin{minipage}[t]{0.48\linewidth}
\small
\begin{tabular}{|l|l|}
\hline
\rowcolor{azulverbo!15} \textbf{Presente} & \\
\hline
yo & escribo \\
\hline
\rowcolor{grisclaro}
tú & escribes \\
\hline
él / ella / Ud. & escribe \\
\hline
\rowcolor{grisclaro}
nosotros/-as & escribimos \\
\hline
vosotros/-as & escribís \\
\hline
\rowcolor{grisclaro}
ellos / ellas / Uds. & escriben \\
\hline
\end{tabular}
\end{minipage}

\vspace{4pt}
\small
\textbf{Frases:}
\begin{itemize}[nosep, leftmargin=1em]
\item Presente: \textit{Elena no \textbf{escribe} en su diario todos los días.}
\item Infinitivo: \textit{A Elena le gusta \textbf{escribir} en su diario.}
\item Participio: \textit{Elena ha \textbf{\textcolor{rojoirr}{escrito}} un correo electrónico.} {\footnotesize (\textcolor{rojoirr}{IRREGULAR})}
\item Gerundio: \textit{Elena está \textbf{escribiendo} en su diario.}
\end{itemize}

\begin{trucoverbo}
\textbf{Regularidad:} Verbo regular en el PRESENTE (-ir, mismas terminaciones que \textit{vivir}). Gerundio regular (-iendo). Atención: el participio es \textcolor{rojoirr}{\textbf{IRREGULAR: escrito}} (no \textit{*escribido}). Es el único verbo de esta sección con participio irregular.
\end{trucoverbo}

\tarjetasep

%%---------- TARJETA 8: LLAMARSE ----------
\verboheader{8}{LLAMARSE}{-ar, 1.\textsuperscript{a} conjugación, regular PRONOMINAL}

\begin{minipage}[t]{0.48\linewidth}
\small
\begin{tabular}{|>{\bfseries}l|l|}
\hline
\rowcolor{azulverbo!15} Forma & \\
\hline
Infinitivo & llamarse \\
\hline
\rowcolor{grisclaro}
Participio & llamado \\
\hline
Gerundio & llamándose \\
\hline
\rowcolor{grisclaro}
Estar + ger. & --- (poco frecuente) \\
\hline
\end{tabular}
\end{minipage}
\hfill
\begin{minipage}[t]{0.48\linewidth}
\small
\begin{tabular}{|l|l|}
\hline
\rowcolor{azulverbo!15} \textbf{Presente} & \\
\hline
yo & me llamo \\
\hline
\rowcolor{grisclaro}
tú & te llamas \\
\hline
él / ella / Ud. & se llama \\
\hline
\rowcolor{grisclaro}
nosotros/-as & nos llamamos \\
\hline
vosotros/-as & os llamáis \\
\hline
\rowcolor{grisclaro}
ellos / ellas / Uds. & se llaman \\
\hline
\end{tabular}
\end{minipage}

\vspace{4pt}
\small
\textbf{Frases:}
\begin{itemize}[nosep, leftmargin=1em]
\item Presente: \textit{Mi hermana \textbf{se llama} Lucía.}
\item Infinitivo: \textit{Es difícil \textbf{llamarse} igual que tu padre.}
\item Participio: \textit{Siempre se ha \textbf{llamado} así.}
\item Gerundio: (poco frecuente con \textit{llamarse} --- se omite ejemplo)
\end{itemize}

\begin{trucoverbo}
\textbf{Regularidad:} El verbo \textit{llamar} es regular (-ar, mismas terminaciones que \textit{estudiar}). La particularidad es que es \textbf{PRONOMINAL}: lleva un pronombre reflexivo (\textbf{me, te, se, nos, os, se}) delante que cambia según la persona. El pronombre va ANTES del verbo conjugado. Conocido desde U01.
\end{trucoverbo}

\newpage

%%---------- TARJETA 9: TENER ----------
\verboheader{9}{TENER}{-er, 2.\textsuperscript{a} conjugación, IRREGULAR}

\begin{minipage}[t]{0.48\linewidth}
\small
\begin{tabular}{|>{\bfseries}l|l|}
\hline
\rowcolor{azulverbo!15} Forma & \\
\hline
Infinitivo & tener \\
\hline
\rowcolor{grisclaro}
Participio & tenido \\
\hline
Gerundio & teniendo \\
\hline
\rowcolor{grisclaro}
Estar + ger. & --- (poco frecuente) \\
\hline
\end{tabular}
\end{minipage}
\hfill
\begin{minipage}[t]{0.48\linewidth}
\small
\begin{tabular}{|l|l|}
\hline
\rowcolor{azulverbo!15} \textbf{Presente} & \\
\hline
yo & \textcolor{rojoirr}{\textbf{tengo}} {\footnotesize (IRREG.)} \\
\hline
\rowcolor{grisclaro}
tú & \textcolor{rojoirr}{\textbf{tienes}} {\footnotesize (e\,$\to$\,ie)} \\
\hline
él / ella / Ud. & \textcolor{rojoirr}{\textbf{tiene}} {\footnotesize (e\,$\to$\,ie)} \\
\hline
\rowcolor{grisclaro}
nosotros/-as & tenemos \\
\hline
vosotros/-as & tenéis \\
\hline
\rowcolor{grisclaro}
ellos / ellas / Uds. & \textcolor{rojoirr}{\textbf{tienen}} {\footnotesize (e\,$\to$\,ie)} \\
\hline
\end{tabular}
\end{minipage}

\vspace{4pt}
\small
\textbf{Frases:}
\begin{itemize}[nosep, leftmargin=1em]
\item Presente: \textit{Jorge \textbf{tiene} dos hermanos.}
\item Infinitivo: \textit{Es genial \textbf{tener} una familia grande.}
\item Participio: \textit{Siempre ha \textbf{tenido} muchos amigos.}
\item Gerundio: (poco frecuente con \textit{tener} --- se omite ejemplo)
\end{itemize}

\begin{trucoverbo}
\textbf{Irregularidad:} Doble irregularidad: (1) \textit{yo} $\to$ \textcolor{rojoirr}{\textbf{tengo}} (raíz ten- + -go, patrón compartido con \textit{venir} $\to$ \textit{vengo}, \textit{poner} $\to$ \textit{pongo}); (2) cambio vocálico \textcolor{rojoirr}{\textbf{e\,$\to$\,ie}} en las personas del singular y 3.\textsuperscript{a} plural (\textit{tienes, tiene, tienen}), pero NO en \textit{nosotros/vosotros} (\textit{tenemos, tenéis}). Participio y gerundio regulares.
\end{trucoverbo}

\tarjetasep

%%---------- TARJETA 10: SER ----------
\verboheader{10}{SER}{-er, IRREGULAR --- repaso U01--U02}

\begin{minipage}[t]{0.48\linewidth}
\small
\begin{tabular}{|>{\bfseries}l|l|}
\hline
\rowcolor{azulverbo!15} Forma & \\
\hline
Infinitivo & ser \\
\hline
\rowcolor{grisclaro}
Participio & sido \\
\hline
Gerundio & siendo \\
\hline
\rowcolor{grisclaro}
Estar + ger. & --- (no se usa con \textit{ser}) \\
\hline
\end{tabular}
\end{minipage}
\hfill
\begin{minipage}[t]{0.48\linewidth}
\small
\begin{tabular}{|l|l|}
\hline
\rowcolor{azulverbo!15} \textbf{Presente} & \\
\hline
yo & \textcolor{rojoirr}{\textbf{soy}} \\
\hline
\rowcolor{grisclaro}
tú & \textcolor{rojoirr}{\textbf{eres}} \\
\hline
él / ella / Ud. & \textcolor{rojoirr}{\textbf{es}} \\
\hline
\rowcolor{grisclaro}
nosotros/-as & \textcolor{rojoirr}{\textbf{somos}} \\
\hline
vosotros/-as & \textcolor{rojoirr}{\textbf{sois}} \\
\hline
\rowcolor{grisclaro}
ellos / ellas / Uds. & \textcolor{rojoirr}{\textbf{son}} \\
\hline
\end{tabular}
\end{minipage}

\vspace{4pt}
\small
\textbf{Frases:}
\begin{itemize}[nosep, leftmargin=1em]
\item Presente: \textit{Mis abuelos \textbf{son} mexicanos.}
\item Infinitivo: \textit{Es divertido \textbf{ser} bilingüe.}
\item Participio: \textit{Siempre ha \textbf{sido} un buen estudiante.}
\item Gerundio: (no se usa con \textit{ser} --- se omite ejemplo)
\end{itemize}

\begin{trucoverbo}
\textbf{Irregularidad:} Verbo completamente irregular --- no sigue ningún patrón de conjugación. Cada forma debe memorizarse. Ya conocido de U01--U02. Se recicla en esta sección en las preguntas con interrogativos (¿Quién es?, ¿Cuál es?, ¿De dónde son?).
\end{trucoverbo}

\newpage

%% ============================================================
%% A-bis. TARJETAS DE CONJUGACIÓN (teal) — 3 tarjetas, 1 por página
%% ============================================================

%%---------- CONJUGACIÓN 1: ENFOQUE -AR ----------
\conjheader{1}{ENFOQUE: GRUPO -AR (1.\textsuperscript{a} conjugación)}

\vspace{6pt}

\begin{center}
\small
\begin{tabular}{|l|>{\columncolor{tealconj!15}}l|l|l|}
\hline
\rowcolor{tealconj!25}
\textbf{Persona} & \textbf{-AR (HABLAR)} & \textbf{-ER (comer)} & \textbf{-IR (vivir)} \\
\hline
yo & habl\textbf{O} & como & vivo \\
\hline
\rowcolor{grisclaro}
tú & habl\textbf{AS} & comes & vives \\
\hline
él / ella / Ud. & habl\textbf{A} & come & vive \\
\hline
\rowcolor{grisclaro}
nosotros/-as & habl\textbf{AMOS} & comemos & vivimos \\
\hline
vosotros/-as & habl\textbf{ÁIS} & coméis & vivís \\
\hline
\rowcolor{grisclaro}
ellos / ellas / Uds. & habl\textbf{AN} & comen & viven \\
\hline
\end{tabular}
\end{center}

\vspace{4pt}
\small

\textbf{Qué comparte -ar con los otros grupos:}
\begin{itemize}[nosep, leftmargin=1.2em]
\item \textbf{Yo} $\to$ los tres terminan en \textbf{-o}. Igual siempre.
\item \textbf{Tú} $\to$ los tres llevan \textbf{-s} al final (-as / -es / -es).
\item \textbf{Ellos} $\to$ los tres añaden \textbf{-n} a la forma de él/ella (-an / -en / -en).
\end{itemize}

\textbf{Qué tiene de diferente -ar:}
\begin{itemize}[nosep, leftmargin=1.2em]
\item La \textbf{vocal protagonista es la A}: -as, -a, -amos, -áis, -an $\to$ 5 de 6 personas llevan A.
\item En \textit{tú, él/ella y ellos} el grupo -ar va por un lado (vocal A) y los grupos -er/-ir van juntos (vocal E).
\item En \textit{nosotros} y \textit{vosotros} los tres grupos son completamente distintos: -\textbf{amos} / -emos / -imos, -\textbf{áis} / -éis / -ís.
\end{itemize}

\textbf{Ejemplos en contexto:}
\begin{itemize}[nosep, leftmargin=1.2em]
\item \textit{Yo hablo español, como fruta y vivo en Madrid.} $\to$ Los tres: -o.
\item \textit{Tú hablas mucho, comes poco y vives cerca.} $\to$ -ar lleva A (\textit{hablas}); los otros, E (\textit{comes, vives}).
\item \textit{Nosotros hablamos, comemos y vivimos aquí.} $\to$ Tres terminaciones distintas: A, E, I.
\end{itemize}

\begin{trucoconjugacion}
\textbf{Truco de la A:} <<El grupo -ar es el grupo de la A. Si ves muchas aes en las terminaciones, es un verbo -ar. Es el grupo más grande: la mayoría de los verbos nuevos que aprendas serán -ar.>>
\end{trucoconjugacion}

\begin{tcolorbox}[colback=tealconj!5, colframe=tealconj!40, boxrule=0.5pt, left=4pt, right=4pt, top=3pt, bottom=3pt]
\textbf{Comparación rápida:} ¿En qué se distingue -ar de los otros? $\to$ En la \textbf{vocal}: donde -er/-ir ponen E, el grupo -ar pone \textbf{A}. Solo en \textit{yo} (-o) coinciden todos.
\end{tcolorbox}

\vspace{4pt}
\textbf{Verbos -ar de nivel A1:}\\
{\footnotesize bailar, buscar, cantar, cenar, cocinar, comprar, desayunar, dibujar, escuchar, estudiar, hablar, lavar, limpiar, llamar(se), llegar, llevar, merendar*, mirar, necesitar, ordenar, pasear, pintar, preguntar, preparar, presentar, terminar, tocar, tomar, trabajar, viajar}\\
{\scriptsize (* merendar tiene cambio vocálico e\,$\to$\,ie, pero en A1 se presenta como vocabulario sin formalizar la irregularidad)}

\newpage

%%---------- CONJUGACIÓN 2: ENFOQUE -ER ----------
\conjheader{2}{ENFOQUE: GRUPO -ER (2.\textsuperscript{a} conjugación)}

\vspace{6pt}

\begin{center}
\small
\begin{tabular}{|l|l|>{\columncolor{tealconj!15}}l|l|}
\hline
\rowcolor{tealconj!25}
\textbf{Persona} & \textbf{-AR (hablar)} & \textbf{-ER (COMER)} & \textbf{-IR (vivir)} \\
\hline
yo & hablo & com\textbf{O} & vivo \\
\hline
\rowcolor{grisclaro}
tú & hablas & com\textbf{ES} & vives \\
\hline
él / ella / Ud. & habla & com\textbf{E} & vive \\
\hline
\rowcolor{grisclaro}
nosotros/-as & hablamos & com\textbf{EMOS} & vivimos \\
\hline
vosotros/-as & habláis & com\textbf{ÉIS} & vivís \\
\hline
\rowcolor{grisclaro}
ellos / ellas / Uds. & hablan & com\textbf{EN} & viven \\
\hline
\end{tabular}
\end{center}

\vspace{4pt}
\small

\textbf{Qué comparte -er con los otros grupos:}
\begin{itemize}[nosep, leftmargin=1.2em]
\item \textbf{Yo} $\to$ los tres terminan en \textbf{-o}. Igual siempre.
\item \textbf{Tú, él/ella, ellos} $\to$ -er y -ir son \textbf{idénticos} en estas 4 personas: -es, -e, -en. Solo -ar es diferente (vocal A).
\item \textbf{Tú} $\to$ los tres llevan \textbf{-s} al final. \textbf{Ellos} $\to$ los tres añaden \textbf{-n}.
\end{itemize}

\textbf{Qué tiene de diferente -er:}
\begin{itemize}[nosep, leftmargin=1.2em]
\item La \textbf{vocal protagonista es la E}: -es, -e, -emos, -éis, -en $\to$ 5 de 6 personas llevan E.
\item Es el grupo más pequeño del español (pocos verbos regulares en -er).
\item Se distingue de -ir \textbf{solo en 2 personas}: nosotros (-\textbf{emos} vs. -imos) y vosotros (-\textbf{éis} vs. -ís). En las otras 4 personas, -er e -ir son gemelos.
\end{itemize}

\textbf{Ejemplos en contexto:}
\begin{itemize}[nosep, leftmargin=1.2em]
\item \textit{Yo como y vivo aquí.} $\to$ -er e -ir iguales: -o.
\item \textit{Tú comes y vives aquí.} $\to$ -er e -ir iguales: -es.
\item \textit{Nosotros comemos pero vivimos lejos.} $\to$ ¡Aquí sí cambian! -Emos vs. -imos.
\end{itemize}

\begin{trucoconjugacion}
\textbf{Truco de la E y el gemelo:} <<El grupo -er es el grupo de la E. Pero cuidado: su gemelo (-ir) copia sus terminaciones en 4 de 6 personas. Solo se separan en \textit{nosotros} y \textit{vosotros}. Si ya sabes -er, casi sabes -ir.>>
\end{trucoconjugacion}

\begin{tcolorbox}[colback=tealconj!5, colframe=tealconj!40, boxrule=0.5pt, left=4pt, right=4pt, top=3pt, bottom=3pt]
\textbf{Comparación rápida:} ¿-er vs. -ar? $\to$ Vocal distinta (E vs. A) en 5 de 6 personas. ¿-er vs. -ir? $\to$ Solo cambian \textit{nosotros} (-emos/-imos) y \textit{vosotros} (-éis/-ís). El resto es idéntico.
\end{tcolorbox}

\vspace{4pt}
\textbf{Verbos -er de nivel A1:}\\
{\footnotesize beber, comer, comprender, correr, creer, leer, responder, vender}

\newpage

%%---------- CONJUGACIÓN 3: ENFOQUE -IR ----------
\conjheader{3}{ENFOQUE: GRUPO -IR (3.\textsuperscript{a} conjugación)}

\vspace{6pt}

\begin{center}
\small
\begin{tabular}{|l|l|l|>{\columncolor{tealconj!15}}l|}
\hline
\rowcolor{tealconj!25}
\textbf{Persona} & \textbf{-AR (hablar)} & \textbf{-ER (comer)} & \textbf{-IR (VIVIR)} \\
\hline
yo & hablo & como & viv\textbf{O} \\
\hline
\rowcolor{grisclaro}
tú & hablas & comes & viv\textbf{ES} \\
\hline
él / ella / Ud. & habla & come & viv\textbf{E} \\
\hline
\rowcolor{grisclaro}
nosotros/-as & hablamos & comemos & viv\textbf{IMOS} \\
\hline
vosotros/-as & habláis & coméis & viv\textbf{ÍS} \\
\hline
\rowcolor{grisclaro}
ellos / ellas / Uds. & hablan & comen & viv\textbf{EN} \\
\hline
\end{tabular}
\end{center}

\vspace{4pt}
\small

\textbf{Qué comparte -ir con los otros grupos:}
\begin{itemize}[nosep, leftmargin=1.2em]
\item \textbf{Yo} $\to$ los tres terminan en \textbf{-o}. Igual siempre.
\item \textbf{Tú, él/ella, ellos} $\to$ -ir usa las \textbf{mismas terminaciones que -er}: -es, -e, -en. Son gemelos en estas 4 personas.
\item \textbf{Tú} $\to$ los tres llevan \textbf{-s} al final. \textbf{Ellos} $\to$ los tres añaden \textbf{-n}.
\end{itemize}

\textbf{Qué tiene de diferente -ir:}
\begin{itemize}[nosep, leftmargin=1.2em]
\item La \textbf{vocal protagonista es la I}, pero \textbf{solo aparece en 2 personas}: nosotros (-\textbf{imos}) y vosotros (-\textbf{ís}). En las otras 4 personas usa la E de -er.
\item Es el único grupo donde la vocal del nombre (-ir) no domina toda la conjugación: la I solo manda en \textit{nosotros/vosotros}.
\end{itemize}

\textbf{Ejemplos en contexto:}
\begin{itemize}[nosep, leftmargin=1.2em]
\item \textit{Él come y vive en Getafe.} $\to$ -er e -ir iguales: -e.
\item \textit{Ellos comen y viven juntos.} $\to$ -er e -ir iguales: -en.
\item \textit{Nosotros comemos, pero vivimos separados.} $\to$ ¡Aquí sí! -emos vs. -\textbf{imos}. La I aparece.
\item \textit{Vosotros coméis y vivís mucho.} $\to$ -éis vs. -\textbf{ís}. La I de nuevo.
\end{itemize}

\begin{trucoconjugacion}
\textbf{Truco del impostor:} <<El grupo -ir es un impostor: se disfraza de -er en 4 de 6 personas. Solo se quita la máscara en \textit{nosotros} (-Imos) y \textit{vosotros} (-Ís). Si dudas entre -er e -ir, piensa: ¿estoy en \textit{nosotros} o \textit{vosotros}? Si no $\to$ da igual, son iguales.>>
\end{trucoconjugacion}

\begin{tcolorbox}[colback=tealconj!5, colframe=tealconj!40, boxrule=0.5pt, left=4pt, right=4pt, top=3pt, bottom=3pt]
\textbf{Comparación rápida:} ¿-ir vs. -er? $\to$ Gemelos en todo menos \textit{nosotros/vosotros}. ¿-ir vs. -ar? $\to$ Solo coinciden en \textit{yo} (-o). En todo lo demás, vocal distinta.
\end{tcolorbox}

\vspace{4pt}
\textbf{Verbos -ir de nivel A1:}\\
{\footnotesize abrir, describir, escribir, recibir, subir, vivir}

\newpage

%% ============================================================
%% B. TARJETAS DE POSESIVOS (morado) — 3 tarjetas, 1 por página
%% ============================================================

%%---------- POSESIVOS 1: MI, TU, SU ----------
\posheader{1}{MI, TU, SU: sin concordancia de género}

\vspace{4pt}
\small
\textit{Estos posesivos NO cambian con el género (masculino/femenino). Solo cambian con el número: ¿una cosa o varias cosas?}

\vspace{6pt}

\begin{center}
\begin{tabular}{|l|l|l|}
\hline
\rowcolor{moradopos!20}
\textbf{Poseedor} & \textbf{1 cosa poseída} & \textbf{Varias cosas poseídas} \\
\hline
\textbf{yo} & \textbf{mi} hermano / \textbf{mi} hermana & \textbf{mis} hermanos / \textbf{mis} hermanas \\
\hline
\rowcolor{grisclaro}
\textbf{tú} & \textbf{tu} primo / \textbf{tu} prima & \textbf{tus} primos / \textbf{tus} primas \\
\hline
\textbf{él / ella / Ud.} & \textbf{su} padre / \textbf{su} madre & \textbf{sus} padres / \textbf{sus} madres \\
\hline
\rowcolor{grisclaro}
\textbf{ellos / ellas / Uds.} & \textbf{su} padre / \textbf{su} madre & \textbf{sus} padres / \textbf{sus} madres \\
\hline
\end{tabular}
\end{center}

\vspace{6pt}

\textbf{¿Qué cambia?}\\
Solo el número: una cosa $\to$ mi, tu, su; varias cosas $\to$ mi\textbf{s}, tu\textbf{s}, su\textbf{s}. Se añade una \textbf{-S}.

\vspace{4pt}
\textbf{¿Qué NO cambia?}\\
El género: \textit{mi hermano} = \textit{mi hermana}. Da igual si es masculino o femenino.

\vspace{4pt}
\textbf{Ojo con SU/SUS:}\\
La misma forma (\textit{su/sus}) sirve para él, ella, usted, ellos, ellas y ustedes. El contexto aclara de quién es: \textit{<<Lucía vive con su madre>>} $\to$ su = de Lucía.

\vspace{4pt}
\textbf{Ya lo conoces de Vocabulario (p.\,34--35):}\\
En los textos sobre las familias de Javier y Lucía ya aparecieron: \textit{su hermano, su abuelo, su madre, su padre, su tía}. Ahora formalizamos la regla completa.

\begin{trucoposesivo}
\textbf{Truco de la S:} <<Mi, tu, su son cortitos y fáciles. No miran si es chico o chica. Solo miran cuántas cosas hay: una cosa $\to$ mi; varias cosas $\to$ mis. Solo añades una S, como en el plural de los nombres.>>
\end{trucoposesivo}

\textbf{Frases de ejemplo:}
\begin{itemize}[nosep, leftmargin=1.2em]
\item \textit{Mi padre se llama Luis y mi madre se llama Alicia.} $\to$ mi + singular, sin cambio de género.
\item \textit{Mis abuelos viven en Sevilla.} $\to$ mis + plural.
\item \textit{¿Tu hermana estudia aquí?} $\to$ tu + singular.
\item \textit{¿Tus primos hablan español?} $\to$ tus + plural.
\item \textit{Su familia es grande.} $\to$ su + singular (¿de quién? $\to$ del contexto).
\item \textit{Sus hermanos comen en el instituto.} $\to$ sus + plural.
\end{itemize}

\newpage

%%---------- POSESIVOS 2: NUESTRO, VUESTRO ----------
\posheader{2}{NUESTRO, VUESTRO: con concordancia de género y número}

\vspace{4pt}
\small
\textit{Estos posesivos SÍ cambian con el género (masculino/femenino) Y con el número (singular/plural). Funcionan como un adjetivo: tienen 4 formas.}

\vspace{6pt}

\begin{center}
\begin{tabular}{|l|l|l|l|l|}
\hline
\rowcolor{moradopos!20}
\textbf{Poseedor} & \textbf{1 cosa (masc.)} & \textbf{1 cosa (fem.)} & \textbf{Varias (masc.)} & \textbf{Varias (fem.)} \\
\hline
\textbf{nosotros/-as} & \textbf{nuestro} padre & \textbf{nuestra} madre & \textbf{nuestros} primos & \textbf{nuestras} primas \\
\hline
\rowcolor{grisclaro}
\textbf{vosotros/-as} & \textbf{vuestro} padre & \textbf{vuestra} madre & \textbf{vuestros} primos & \textbf{vuestras} primas \\
\hline
\end{tabular}
\end{center}

\vspace{6pt}

\textbf{¿Qué cambia?}\\
Todo: género \textbf{y} número. Cuatro terminaciones como un adjetivo:

\vspace{4pt}
\begin{center}
\begin{tabular}{|l|c|c|}
\hline
\rowcolor{moradopos!15}
 & \textbf{Singular} & \textbf{Plural} \\
\hline
\textbf{Masculino} & nuestr\textbf{o} & nuestr\textbf{os} \\
\hline
\rowcolor{grisclaro}
\textbf{Femenino} & nuestr\textbf{a} & nuestr\textbf{as} \\
\hline
\end{tabular}
\end{center}

\vspace{2pt}
{\footnotesize (Igual para vuestr-o / vuestr-a / vuestr-os / vuestr-as.)}

\vspace{6pt}

\textbf{Comparación con la tarjeta 1:}
\begin{itemize}[nosep, leftmargin=1.2em]
\item \textit{Mi, tu, su} $\to$ solo cambian singular/plural (+S). Fáciles.
\item \textit{Nuestro, vuestro} $\to$ cambian como adjetivos (-o/-a/-os/-as). Hay que fijarse en la cosa poseída.
\end{itemize}

\begin{trucoposesivo}
\textbf{Truco del adjetivo:} <<Nuestro y vuestro se comportan como adjetivos: miran la cosa poseída y copian su género y número. ¿La cosa es femenina? $\to$ nuestr-A. ¿Son varias cosas femeninas? $\to$ nuestr-AS. Piénsalo como \textit{bonito/bonita/bonitos/bonitas}: misma lógica.>>
\end{trucoposesivo}

\textbf{Frases de ejemplo:}
\begin{itemize}[nosep, leftmargin=1.2em]
\item \textit{Nuestro padre trabaja en Madrid.} $\to$ masc. singular.
\item \textit{Nuestra madre habla tres idiomas.} $\to$ fem. singular.
\item \textit{Nuestros abuelos viven en Getafe.} $\to$ masc. plural.
\item \textit{Nuestras hermanas estudian música.} $\to$ fem. plural.
\item \textit{¿Vuestro instituto es grande?} $\to$ masc. singular.
\item \textit{¿Vuestra profesora es española?} $\to$ fem. singular.
\item \textit{¿Vuestros primos comen aquí?} $\to$ masc. plural.
\item \textit{¿Vuestras amigas viven cerca?} $\to$ fem. plural.
\end{itemize}

\newpage

%%---------- POSESIVOS 3: RESUMEN ----------
\posheader{3}{Resumen y estrategia global}

\vspace{6pt}
\small

\begin{tcolorbox}[colback=moradopos!6, colframe=moradopos!50, boxrule=0.8pt, left=6pt, right=6pt, top=6pt, bottom=6pt, title={\textbf{Dos preguntas para elegir la forma correcta}}, coltitle=white, fonttitle=\bfseries]
\textbf{Pregunta 1: ¿De quién es?} $\to$ Elige el posesivo (mi / tu / su / nuestro / vuestro).

\vspace{4pt}
\textbf{Pregunta 2: ¿Cuántas cosas son (y de qué género)?}
\begin{itemize}[nosep, leftmargin=1.2em]
\item Si es \textit{mi, tu, su} $\to$ solo mira el número: una cosa = mi; varias = mis (+S).
\item Si es \textit{nuestro, vuestro} $\to$ mira género y número: cambia como un adjetivo (-o/-a/-os/-as).
\end{itemize}
\end{tcolorbox}

\vspace{8pt}

\textbf{Mapa visual rápido:}

\vspace{4pt}
\begin{center}
\begin{tabular}{|l|l|l|}
\hline
\rowcolor{moradopos!20}
 & \textbf{Solo cambia número (+S)} & \textbf{Cambia género y número (-o/-a/-os/-as)} \\
\hline
\textbf{1 dueño} & mi/mis, tu/tus, su/sus & --- \\
\hline
\rowcolor{grisclaro}
\textbf{Varios dueños} & su/sus & nuestro/-a/-os/-as, vuestro/-a/-os/-as \\
\hline
\end{tabular}
\end{center}

\vspace{10pt}

\begin{trucoposesivo}
\textbf{Truco final:} <<Los posesivos cortos (mi, tu, su) son fáciles: +S y listo. Los posesivos largos (nuestro, vuestro) son como adjetivos. Y \textit{su/sus} es el comodín: sirve para un dueño o varios.>>
\end{trucoposesivo}

\newpage

%% ============================================================
%% C. TARJETAS DE INTERROGATIVOS (naranja) — 3 tarjetas, 1 por página
%% ============================================================

%%---------- INTERROGATIVOS 1: INVARIABLES ----------
\interrheader{1}{INVARIABLES: no cambian nunca}

\vspace{4pt}
\small
\textit{Estas formas son siempre iguales. No importa si la respuesta es masculina, femenina, singular o plural: el interrogativo no cambia.}

\vspace{6pt}

\begin{center}
\begin{tabular}{|l|p{3.2cm}|l|p{5.5cm}|}
\hline
\rowcolor{naranjainterr!20}
\textbf{Interrogativo} & \textbf{Pregunta por...} & \textbf{¿Con qué va?} & \textbf{Ejemplo} \\
\hline
\textbf{¿Quién?} & una persona (identidad) & + verbo & \textit{¿Quién es tu hermano?} --- Nacho. \\
\hline
\rowcolor{grisclaro}
\textbf{¿Qué?} & una cosa / definición & + verbo & \textit{¿Qué estudia Lucía?} --- Música. \\
\hline
\textbf{¿Cómo?} & descripción / manera & + verbo & \textit{¿Cómo se llama tu madre?} --- Alicia. \\
\hline
\rowcolor{grisclaro}
\textbf{¿Dónde?} & un lugar & + verbo & \textit{¿Dónde vive Javier?} --- En Getafe. \\
\hline
\textbf{¿De dónde?} & un origen / procedencia & + verbo & \textit{¿De dónde es tu padre?} --- De México. \\
\hline
\end{tabular}
\end{center}

\vspace{6pt}

\textbf{¿Qué tienen en común?}
\begin{itemize}[nosep, leftmargin=1.2em]
\item Los 5 van seguidos de un \textbf{verbo} (nunca de un sustantivo directamente).
\item Los 5 son \textbf{invariables}: no cambian ni en género ni en número. Se memorizan una vez y sirven siempre.
\item Todos llevan \textbf{tilde} (quién, qué, cómo, dónde). Sin tilde son otra cosa (\textit{que, como, donde} $\to$ no son preguntas).
\end{itemize}

\begin{trucointerr}
\textbf{Truco de los 5 fijos:} <<Quién, qué, cómo, dónde, de dónde: cinco palabras que nunca cambian. Siempre van con un verbo detrás. Apréndelas una vez y ya está.>>
\end{trucointerr}

\textbf{Más frases de ejemplo:}
\begin{itemize}[nosep, leftmargin=1.2em]
\item \textit{¿Quién vive en Buenos Aires?} --- Su hermano.
\item \textit{¿Qué come Jorge?} --- Fruta y pasta.
\item \textit{¿Cómo es tu familia?} --- Grande y divertida.
\item \textit{¿Dónde trabaja tu padre?} --- En un hospital.
\item \textit{¿De dónde son tus abuelos?} --- De Sevilla.
\end{itemize}

\newpage

%%---------- INTERROGATIVOS 2: CON CONCORDANCIA ----------
\interrheader{2}{CON CONCORDANCIA: cambian según el sustantivo}

\vspace{4pt}
\small
\textit{Estas formas SÍ cambian. Cuánto mira el género y el número del sustantivo que acompaña. Cuál solo mira el número.}

\vspace{8pt}

\noindent\colorbox{naranjainterr!15}{\parbox{\dimexpr\linewidth-2\fboxsep}{\textbf{\large CUÁNTO/-A/-OS/-AS} $\to$ pregunta por cantidad}}

\vspace{4pt}

\begin{minipage}[t]{0.45\linewidth}
\begin{center}
\begin{tabular}{|l|c|c|}
\hline
\rowcolor{naranjainterr!15}
 & \textbf{Singular} & \textbf{Plural} \\
\hline
\textbf{Masculino} & cuánt\textbf{o} & cuánt\textbf{os} \\
\hline
\rowcolor{grisclaro}
\textbf{Femenino} & cuánt\textbf{a} & cuánt\textbf{as} \\
\hline
\end{tabular}
\end{center}

\vspace{4pt}
\begin{itemize}[nosep, leftmargin=1em, font=\small]
\item ¿Con qué va? $\to$ + \textbf{sustantivo} (la cosa que se cuenta).
\item Cambia como un adjetivo: mira el sustantivo y copia su género y número.
\item Conexión con posesivos: igual que \textit{nuestro/-a/-os/-as}, tiene 4 formas.
\end{itemize}
\end{minipage}
\hfill
\begin{minipage}[t]{0.52\linewidth}
\begin{tabular}{|p{5cm}|l|}
\hline
\rowcolor{naranjainterr!10}
\textbf{Ejemplo} & \textbf{Gén./núm.} \\
\hline
\textit{¿\textbf{Cuántos} hermanos tienes?} & masc. plural \\
\hline
\rowcolor{grisclaro}
\textit{¿\textbf{Cuántas} hermanas tienes?} & fem. plural \\
\hline
\textit{¿\textbf{Cuánto} tiempo estudias?} & masc. singular \\
\hline
\rowcolor{grisclaro}
\textit{¿\textbf{Cuánta} fruta comes?} & fem. singular \\
\hline
\end{tabular}
\end{minipage}

\vspace{10pt}

\noindent\colorbox{naranjainterr!15}{\parbox{\dimexpr\linewidth-2\fboxsep}{\textbf{\large CUÁL / CUÁLES} $\to$ pregunta por elección (elegir entre opciones)}}

\vspace{4pt}

\begin{minipage}[t]{0.45\linewidth}
\begin{center}
\begin{tabular}{|c|c|}
\hline
\rowcolor{naranjainterr!15}
\textbf{Singular} & \textbf{Plural} \\
\hline
cuál & cuáles \\
\hline
\end{tabular}
\end{center}

\vspace{4pt}
\begin{itemize}[nosep, leftmargin=1em, font=\small]
\item ¿Con qué va? $\to$ + verbo (generalmente \textit{ser}). También puede ir + \textbf{de} + sustantivo.
\item Solo cambia en número (singular/plural). No cambia en género.
\item Conexión con posesivos: como \textit{mi/mis} $\to$ solo se añade una terminación para plural.
\end{itemize}
\end{minipage}
\hfill
\begin{minipage}[t]{0.52\linewidth}
\begin{tabular}{|p{5.5cm}|l|}
\hline
\rowcolor{naranjainterr!10}
\textbf{Ejemplo} & \textbf{Número} \\
\hline
\textit{¿\textbf{Cuál} es tu nombre?} & singular \\
\hline
\rowcolor{grisclaro}
\textit{¿\textbf{Cuáles} son tus asignaturas favoritas?} & plural \\
\hline
\textit{¿\textbf{Cuál} de tus hermanos es mayor?} & singular \\
\hline
\end{tabular}
\end{minipage}

\vspace{8pt}

\textbf{Comparación entre los dos:}
\begin{itemize}[nosep, leftmargin=1.2em]
\item \textit{Cuánto} $\to$ cambia \textbf{género y número} (4 formas) $\to$ va con \textbf{sustantivo}.
\item \textit{Cuál} $\to$ cambia solo \textbf{número} (2 formas) $\to$ va con \textbf{verbo} (o \textit{de} + sustantivo).
\end{itemize}

\begin{trucointerr}
\textbf{Truco del adjetivo y el pronombre:} <<Cuánto funciona como un adjetivo: acompaña a un sustantivo y copia su género y número (\textit{cuántos hermanos, cuántas hermanas}). Cuál funciona como un pronombre: va con verbo y solo cambia singular/plural (\textit{cuál es, cuáles son}).>>
\end{trucointerr}

\newpage

%%---------- INTERROGATIVOS 3: RESUMEN ----------
\interrheader{3}{Resumen y estrategia global}

\vspace{6pt}
\small

\begin{tcolorbox}[colback=naranjainterr!6, colframe=naranjainterr!50, boxrule=0.8pt, left=6pt, right=6pt, top=6pt, bottom=6pt, title={\textbf{Tres preguntas para elegir el interrogativo correcto}}, coltitle=white, fonttitle=\bfseries]
\textbf{Pregunta 1: ¿Qué quiero saber?}
\begin{itemize}[nosep, leftmargin=1.2em]
\item Una persona $\to$ ¿Quién?
\item Una cosa / definición $\to$ ¿Qué?
\item Una descripción / manera $\to$ ¿Cómo?
\item Un lugar $\to$ ¿Dónde? / ¿De dónde?
\item Una cantidad $\to$ ¿Cuánto/-a/-os/-as?
\item Una elección $\to$ ¿Cuál / Cuáles?
\end{itemize}

\vspace{4pt}
\textbf{Pregunta 2: ¿Cambia o no cambia?}
\begin{itemize}[nosep, leftmargin=1.2em]
\item Quién, qué, cómo, dónde, de dónde $\to$ \textbf{NO cambian nunca}.
\item Cuánto $\to$ cambia género y número (4 formas, como un adjetivo).
\item Cuál $\to$ cambia solo número (2 formas).
\end{itemize}

\vspace{4pt}
\textbf{Pregunta 3: ¿Va con verbo o con sustantivo?}
\begin{itemize}[nosep, leftmargin=1.2em]
\item Con \textbf{verbo}: quién, qué, cómo, dónde, de dónde, cuál.
\item Con \textbf{sustantivo}: cuánto/-a/-os/-as.
\end{itemize}
\end{tcolorbox}

\vspace{10pt}

\textbf{Mapa visual rápido:}

\vspace{4pt}
\begin{center}
\begin{tabular}{|l|c|l|l|}
\hline
\rowcolor{naranjainterr!20}
\textbf{Interrogativo} & \textbf{¿Cambia?} & \textbf{¿Va con...?} & \textbf{Pregunta por...} \\
\hline
¿Quién? & No & verbo & persona \\
\hline
\rowcolor{grisclaro}
¿Qué? & No & verbo & cosa \\
\hline
¿Cómo? & No & verbo & descripción \\
\hline
\rowcolor{grisclaro}
¿Dónde? & No & verbo & lugar \\
\hline
¿De dónde? & No & verbo & origen \\
\hline
\rowcolor{grisclaro}
¿Cuánto/-a/-os/-as? & Género + número & sustantivo & cantidad \\
\hline
¿Cuál / Cuáles? & Número & verbo & elección \\
\hline
\end{tabular}
\end{center}

\vspace{10pt}

\begin{trucointerr}
\textbf{Truco final:} <<5 fijos + 2 que cambian. Los 5 fijos van siempre con verbo. De los 2 que cambian: \textit{cuánto} es un adjetivo (va con sustantivo, 4 formas) y \textit{cuál} es un pronombre (va con verbo, 2 formas). Todos llevan tilde: si no hay tilde, no es pregunta.>>
\end{trucointerr}

\end{document}
